% Part: proof-theory
% Chapter: cut-elimination
% Section: ce-topmost

\documentclass[../../../include/open-logic-section]{subfiles}

\begin{document}

\olfileid{pt}{cut}{top}

\olsection{सर्वोच्च स्थानी असलेले कट काढणे}

\begin{defn}
एका \CutCS{} नियमाने संपणारी !!a{proof}~$\pi$
विचारात घ्या:
\[
\AxiomC{}
\RightLabel{$\pi_1$}
\Deduce$\Gamma \fCenter \Delta, !A$
\AxiomC{}
\RightLabel{$\pi_2$}
\Deduce$!A, \Gamma \fCenter \Delta$
\RightLabel{\CutCS}
\BinaryInf$\Gamma \fCenter \Delta$
\DisplayProof
\]
तिच्यात अन्य \CutCS{} अनुमाने नसावीत.
$\pi$ ची \emph{कट-उंची} म्हणजे
$\cheight{\pi} = \pheight{\pi_1} + \pheight{\pi_2}$.
$\pi$ ची \emph{कट-कोटी} म्हणजे
$\cutrank{\pi} = \depth{!A}$.
\end{defn}

\begin{lem}\ollabel{lem:cut-adm-G3c}
\CutCS{} नियम~\Log{G3c} मध्ये ग्राह्य आहे.
दुसऱ्या शब्दांत, !!a{proof}~$\pi$ ही एका
\CutCS{} ने संपत असेल आणि इतरत्र
फक्त~$\Log{G3c}$ चे नियम वापरत असेल,
तर त्याच अंतिम क्रमवर्तीची
\Log{G3c} मधील !!a{proof}~$\pi'$ असते.
\end{lem}

\begin{proof}
कटची !!{height} आणि कोटी यांवरील
दुहेरी विगमनाने सिद्धता करू.
!!{proof}~$\pi$ चा शेवट असा आहे:
\[
\AxiomC{}
\RightLabel{$\pi_1$}
\Deduce$\Gamma \fCenter \Delta, !A$
\AxiomC{}
\RightLabel{$\pi_2$}
\Deduce$!A, \Gamma \fCenter \Delta$
\RightLabel{\CutCS}
\BinaryInf$\Gamma \fCenter \Delta$
\DisplayProof
\]

विगमनाचा आधार: $\cheight{\pi} = 0$
आणि $\cutrank{\pi} = 0$.
$\cheight{\pi} = \pheight{\pi_1}
+ \pheight{\pi_2} = 0$ असल्याने
$\Gamma \Sequent \Delta, !A$ आणि
$!A, \Gamma \Sequent \Delta$
ही दोन्ही स्वयंसिद्धके आहेत.
दोन प्रसंग आहेत: (a)~$!A$ हा
त्यांपैकी एखाद्यात मुख्य नाही,
किंवा (b)~$!A$ हा दोन्हींत मुख्य आहे.
खाली दोन्ही प्रसंगांत
$\Gamma \Sequent \Delta$ हे स्वयंसिद्धक
असलेच पाहिजे हे दाखवू
(त्यासाठी विगमन गृहीतक लागत नाही).

आता $\pi_1$ आणि $\pi_2$ स्वयंसिद्धके
आहेत की अनुमानाने संपतात, तसेच त्या
अनुमानात~$!A$ मुख्य आहे की नाही,
यावरून प्रसंग वेगळे करू. प्रसंग A व~B
विगमनाचा आधार समाविष्ट करतात.
विगमन पायरीत $\cheight{\pi} > 0$
किंवा $\cutrank{\pi} > 0$
(किंवा दोन्ही) असे मानतो. विगमन
गृहीतकानुसार, एका \CutCS{} ने
संपणाऱ्या !!{proof} मधील अंतिम
\CutCS{} काढता येतो, जर तिची
कट-कोटी $= \cutrank{\pi}$ आणि
कटची !!{height} $< \cheight{\pi}$
असेल, किंवा तिची कट-कोटी
$< \cutrank{\pi}$ असेल.
$\cheight{\pi} = 0$ चा प्रसंग
(दोन्ही आधार स्वयंसिद्धके) A व~B
मध्ये येतो. $\cheight{\pi} > 0$
चे प्रसंग C--D मध्ये येतात.

\paragraph{प्रसंग A:}
काही $!B$ हे $\Gamma$ आणि~$\Delta$
या दोन्हींत येते. $!B$ अणुसूत्र
असेल, तर $\Gamma \Sequent \Delta$
हे स्वयंसिद्धक असते; अन्यथा
\olref[seq][ex]{prop:G3c-axioms} नुसार
त्याची \Log{G3c} मधील \CutCS{}
विरहित !!a{proof}~$\pi'$ असते.
यात विगमन आधाराचा प्रसंग~(a) येतो:
$\Gamma \Sequent \Delta, !A$
किंवा $!A, \Gamma \Sequent \Delta$
हे स्वयंसिद्धक असून त्यात $!A$
मुख्य नसेल, तर अन्य अणुसूत्र~$!B$
हे $\Gamma$ आणि~$\Delta$ या
दोन्हींत असते. विगमन पायरीतही
एखादा आधार स्वयंसिद्धक असून त्यात
$!A$ मुख्य नसेल, तर हा प्रसंग
लागू होतो; दुसरा आधार नियमाचा
निष्कर्ष असला तरी चालते.

\paragraph{प्रसंग B:}
दोन्ही आधार स्वयंसिद्धके आहेत
आणि $!A$ मुख्य, म्हणून अणुसूत्र,
आहे. डावा आधार
$\Gamma \Sequent \Delta, !A$
पाहा. या $\Gamma \Sequent \Delta, !A$
मध्ये $!A$ मुख्य असल्यामुळे ते
$\Gamma$ मध्ये असलेच पाहिजे
(नसते तर ही क्रमवर्ती स्वयंसिद्धक
ठरली नसती). त्याचप्रमाणे,
$!A$ हे $\Delta$ मध्येही असलेच
पाहिजे; अन्यथा
$!A, \Gamma \Sequent \Delta$
स्वयंसिद्धक ठरणार नाही.
त्यामुळे $!A$ अणुसूत्र असून
$\Gamma$ आणि $\Delta$ या
दोन्हींत येते; म्हणजे
$\Gamma \Sequent \Delta$
हे स्वयंसिद्धक आहे. हा
विगमन आधाराचा प्रसंग~(b) आहे.

\begin{center}
प्रसंग A आणि B चे पर्यायी रूप
\end{center}

\paragraph{प्रसंग A:}
कटच्या आधारांपैकी एक
$!C, \Gamma' \Sequent \Delta', !C$
या रूपाचे स्वयंसिद्धक आहे,
जेथे $!C$ अणुसूत्र आहे.
$!C$ हे कटचे !!{formula}~$!A$
नसेल, तर $!C$ हे $\Gamma$
आणि~$\Delta$ या दोन्हींत
असलेच पाहिजे. मग
$\Gamma \Sequent \Delta$
हे स्वयंसिद्धक असून तेच~$\pi'$
घेता येते. $!C \ident !A$
कटचे !!{formula} असेल, तर
डावा आधार स्वयंसिद्धक असल्यास
$!A \in \Gamma$ आणि
$\Gamma = \Gamma', !A$;
उजवा आधार स्वयंसिद्धक असल्यास
$!A \in \Delta$ आणि
$\Delta = \Delta', !A$.
पहिल्या प्रसंगात $\pi_2$
ही $!A, !A, \Gamma' \Sequent \Delta$
ची !!a{proof} असते;
\LeftR{\Contraction} ची ग्राह्यता
(\olref[seq][inv]{prop:G3c-cont-adm})
वापरून $\Gamma \Sequent \Delta$
ची !!a{proof} मिळते.
दुसऱ्या प्रसंगात $\pi_1$
ही $\Gamma \Sequent \Delta', !A, !A$
ची !!a{proof} असते;
\RightR{\Contraction} ची ग्राह्यता
वापरून $\Gamma \Sequent \Delta$
ची !!a{proof} मिळते.
%READERNOTE{OLINC-341: स्रोत दुसऱ्या प्रसंगात pi_2 लिहितो; उजवा आधार स्वयंसिद्धक असल्यास आवश्यक सिद्धता डाव्या आधाराची pi_1 असते.}

\paragraph{प्रसंग B:}
एखादा आधार
$\lfalse, \Gamma' \Sequent \Delta'$
या रूपाचा स्वयंसिद्धक आहे.
डावा आधार असा असेल, तर
$\lfalse \in \Gamma$ आणि
$\Gamma \Sequent \Delta$
हेही स्वयंसिद्धक आहे. उजवा
आधार असा असेल, तर
$\lfalse \in \Gamma$
(मग पुन्हा $\Gamma \Sequent \Delta$
स्वयंसिद्धक असते) किंवा
$!A \ident \lfalse$.
दुसऱ्या उपप्रसंगात $\pi_1$
ही $\Gamma \Sequent \Delta, \lfalse$
ची !!a{proof} आहे. $\pi_1$
मधील प्रत्येक क्रमवर्तीच्या
उत्तरांगातून $\lfalse$ ची
एक प्रत काढून
$\Gamma \Sequent \Delta$
ची !!a{proof} मिळते.
(सराव: $\pheight{\pi_1}$
वरील विगमनाने हे तपासा.)

आता प्रसंग~A किंवा B लागू
होत नाही असे समजा. उर्वरित
प्रसंगांत $\cheight{\pi} > 0$;
म्हणजे किमान एका आधाराचा
शेवट नियमाच्या अनुमानाने होतो.

\paragraph{प्रसंग C आणि D: कटचे स्थान बदलणे.}
प्रसंग~C आणि~D मधील सर्व शक्यतांत
एक सामाईक नमुना आहे. सुरुवातीला
\CutCS{} ने संपणारी !!a{proof} असते;
तिच्या एका आधाराचा निष्कर्ष नियम~$R$
ने मिळतो, पण कटचे !!{formula}~$!A$
त्या नियमात मुख्य नसते (काही अन्य
!!{formula}~$!D$ मुख्य असते).
या !!{proof} मध्ये \CutCS{}
हा नियम~$R$ नंतर येतो; तिचे
योजनात्मक रूप असे:
\[
\AxiomC{}
\RightLabel{$\pi_1$}
\Deduce$\Gamma \fCenter \Delta', !D, !A$
\AxiomC{}
\RightLabel{$\pi_3$}
\Deduce$!A, \Gamma, \Pi \fCenter \Delta', \Lambda$
\RightLabel{$R$}
\UnaryInf$!A, \Gamma \fCenter \Delta', !D$
\RightSubproofLabel{$\pi_2$}
\RightLabel{\CutCS}
\BinaryInf$\Gamma \fCenter \Delta', !D$
\DisplayProof
\]
हे रूप फक्त त्या प्रसंगाचे आहे
जेथे $R$ हा एक-आधाराचा उजवा
नियम आहे, $A$ हे $R$ चे
मुख्य !!{formula} नाही, आणि
\emph{मुख्य} असलेले !!{formula}~$!D$
हे \CutCS{} च्या उजव्या
आधाराच्या उत्तरांगात आहे.
$!D$ पूर्वांगात असेल (म्हणजे
$R$ डावा नियम असेल) किंवा
\CutCS{} च्या डाव्या आधारात
असेल, तर रूप अर्थातच बदलेल.
$\Pi$ आणि~$\Lambda$ मधील
!!{formula}s या नियम~$R$
च्या सक्रिय !!{formula}s दर्शवतात.

आता हा क्रम उलटून
$R$ हा~\CutCS{} नंतर येणारी
!!a{proof} विचारात घ्या:
\[
\AxiomC{}
\RightLabel{$\pi_1'$}
\Deduce$\Gamma, \Pi \fCenter \Delta', !D, \Lambda, !A$
\AxiomC{}
\RightLabel{$\pi_3'$}
\Deduce$!A, \Gamma, \Pi \fCenter \Delta', !D, \Lambda$
\RightLabel{\CutCS}
\BinaryInf$\Gamma, \Pi \fCenter \Delta', !D, \Lambda$
\RightLabel{$R$}
\UnaryInf$\Gamma, \fCenter \Delta', !D, !D$
\DisplayProof
\]
$\pi_1$ आणि $\pi_3$ यांवर
!!{height} न वाढवणारे दुर्बलीकरण
लावून अनुक्रमे $\pi_1'$ आणि
$\pi_3'$ मिळवतो.

\CutCS{} आणि~$R$ यांचा क्रम
बदलत असल्यामुळे याला
“कटला~$R$ च्या वर नेणे” म्हणतात.
याने \CutCS{} च्या उजव्या
आधारावर संपणाऱ्या !!{proof}
ची उंची कमी होते; म्हणून कटची
!!{height} कमी होते. म्हणजे
नव्या \CutCS{} वरील
उप-!!{proof}s ची !!{height}
$\pi_1$ आणि $\pi_3$ यांच्या
उंचीपेक्षा मोठी नसते, म्हणून
कटची !!{height} कमी होते
(उजवीकडील एक अनुमान काढले आहे).
आता विगमन गृहीतक लावून \CutCS{}
अनुमान काढता येते. शेवटी
$\RightR{\Contraction}$ ची ग्राह्यता
वापरून~$!D$ ची अतिरिक्त
प्रत काढतो.

\paragraph{प्रसंग C:}
$\pi_1$ आणि $\pi_2$ यांपैकी
किमान एकाचा शेवट एक-आधाराच्या
नियम~$R$ ने होतो, आणि त्यात
$!A$ मुख्य नसते. एकूण 18 प्रसंग
आहेत: \CutCS{} च्या डाव्या की
उजव्या आधारात $!A$ मुख्य नाही
आणि नऊ एक-आधाराच्या नियमांपैकी
(\LeftR{\lnot}, \RightR{\lnot},
\RightR{\lor}, \LeftR{\land},
\RightR{\lif}, आणि चार संख्यक नियम)
$R$ कोणता आहे, यावर ते अवलंबून
असतात. आपण फक्त दोन उदाहरणे
पाहू: $\pi_2$ च्या अंतिम
अनुमानात $!A$ मुख्य नाही आणि ते
अनुमान $\RightR{\lif}$ किंवा
\LeftR{\lexists} आहे.
इतर प्रसंग यांसारखेच असून
सरावासाठी ठेवतो.

समजा $\pi$ चा शेवट असा आहे:
\[
\AxiomC{}
\RightLabel{$\pi_1$}
\Deduce$\Gamma \fCenter \Delta', !B \lif !C, !A$
\AxiomC{}
\RightLabel{$\pi_3$}
\Deduce$!A, !B, \Gamma \fCenter \Delta', !C$
\RightLabel{$\RightR{\lif}$}
\UnaryInf$!A, \Gamma \fCenter \Delta', !B \lif !C$
\RightSubproofLabel{$\pi_2$}
\RightLabel{\CutCS}
\BinaryInf$\Gamma \fCenter \Delta', !B \lif !C$
\DisplayProof
\]
येथे $\Delta = \Delta', !B \lif !C$.
$\pi_1$ ला !!{height} न वाढवणारे
दुर्बलीकरण
(\olref[seq][adm]{prop:weak-G3c-adm})
लावून $!B, \Gamma \fCenter
\Delta', !B \lif !C, !C, !A$
ची !!a{proof} $\pi_1'$ मिळवतो
(डावीकडे $!B$ आणि उजवीकडे
$!C$ घालतो). त्याचप्रमाणे
\olref[seq][adm]{prop:weak-G3c-adm}
हे~$\pi_3$ ला लावून
$!A, !B, \Gamma \fCenter
\Delta', !B \lif !C, !C$
ची !!a{proof}~$\pi_3'$ मिळते
(उजवीकडे $!B \lif !C$ घालतो).
विगमन गृहीतक पुढील सिद्धतेला
लागू होते:
\[
\AxiomC{}
\RightLabel{$\pi_1'$}
\Deduce$!B, \Gamma \fCenter \Delta', !B \lif !C, !C, !A$
\AxiomC{}
\RightLabel{$\pi_3'$}
\Deduce$!A, !B, \Gamma \fCenter \Delta', !B \lif !C, !C$
\RightLabel{\CutCS}
\BinaryInf$!B, \Gamma \fCenter \Delta', !B \lif !C, !C$
\DisplayProof
\]
यातील कटचे !!{formula}~$!A$
आहे: या !!{proof} ची कट-कोटी
$\pi$ इतकीच, पण कटची !!{height}
कमी आहे. त्यामुळे \CutCS{}
विरहित \Log{G3c} मधील
$!B, \Gamma \fCenter
\Delta', !B \lif !C, !C$
ची !!a{proof}~$\pi_4$ मिळते.
$\RightR{\lif}$ लावून
$\Gamma \fCenter \Delta',
!B \lif !C, !B \lif !C$
ची !!a{proof} मिळते:
\[
\AxiomC{}
\RightLabel{$\pi_4$}
\Deduce$!B, \Gamma \fCenter \Delta', !B \lif !C, !C$
\RightLabel{$\RightR{\lif}$}
\UnaryInf$\Gamma \fCenter \Delta', !B \lif !C, !B \lif !C$
\DisplayProof
\]
$\Gamma \Sequent \Delta$,
म्हणजे $\Gamma \fCenter
\Delta', !B \lif !C$,
याची !!a{proof}~$\pi'$
मिळवण्यासाठी
\olref[seq][inv]{prop:G3c-cont-adm},
म्हणजे आकुंचनाची ग्राह्यता,
वापरा.

आता $\pi$ चा शेवट
पुढीलप्रमाणे आहे असे समजा:
\[
\AxiomC{}
\RightLabel{$\pi_1$}
\Deduce$\lexists[x][!B(x)], \Gamma' \fCenter \Delta, !A$
\AxiomC{}
\RightLabel{$\pi_3$}
\Deduce$!A, !B(c), \Gamma' \fCenter \Delta$
\RightLabel{$\LeftR{\lexists}$}
\UnaryInf$!A, \lexists[x][!B(x)], \Gamma' \fCenter \Delta$
\RightSubproofLabel{$\pi_2$}
\RightLabel{\CutCS}
\BinaryInf$\Gamma' \fCenter \Delta$
\DisplayProof
\]
पुन्हा पुढील सिद्धतेला
विगमन गृहीतक लावू:
\[
\AxiomC{}
\RightLabel{$\pi_1[!B(c) \Sequent {}]$}
\Deduce$\lexists[x][!B(x)], !B(c), \Gamma' \fCenter \Delta, !A$
\AxiomC{}
\LeftLabel{$\pi_3[\lexists[x][!B(x)] \Sequent {}]$}
\Deduce$!A, \lexists[x][!B(x)], !B(c), \Gamma' \fCenter \Delta$
\RightLabel{\CutCS}
\BinaryInf$\lexists[x][!B(x)], !B(c), \Gamma' \fCenter \Delta$
\RightLabel{\LeftR{\lexists}}
\UnaryInf$\lexists[x][!B(x)], \lexists[x][!B(x)], \Gamma' \fCenter \Delta$
\DisplayProof
\]
येथे आयगेन चराची अट पूर्ण होते.
कारण $\pi_2$ मधील
\LeftR{\lexists} वरील
आयगेन चराच्या अटीनुसार
$c$ हे $\lexists[x][!B(x)]$,
$\Gamma'$ किंवा~$\Delta$
यांत येत नाही. शेवटी
\olref[seq][inv]{prop:G3c-cont-adm}
लावा.

\paragraph{प्रसंग D:}
$\pi_1$ आणि $\pi_2$ यांपैकी
किमान एकाचा शेवट दोन-आधारांच्या
नियम~$R$ ने होतो आणि त्यात
$!A$ मुख्य नसते. असे सहा
प्रसंग आहेत. $\pi_2$ चा
शेवट~$\RightR{\land}$ ने होत
असताना त्याच्या अंतिम अनुमानात
$!A$ मुख्य नसतो हा प्रसंग
पाहू; इतर तसेच आहेत.

समजा $\pi$ चा शेवट असा आहे:
\[
\AxiomC{}
\RightLabel{$\pi_1$}
\Deduce$\Gamma \fCenter \Delta', !B \land !C, !A$
\AxiomC{}
\RightLabel{$\pi_3$}
\Deduce$!A, \Gamma \fCenter \Delta', !B$
\AxiomC{}
\RightLabel{$\pi_4$}
\Deduce$!A, \Gamma \fCenter \Delta', !C$
\RightLabel{$\RightR{\land}$}
\BinaryInf$!A, \Gamma \fCenter \Delta', !B \land !C$
\RightSubproofLabel{$\pi_2$}
\RightLabel{\CutCS}
\BinaryInf$\Gamma \fCenter \Delta$
\DisplayProof
\]
विगमन गृहीतक पुढील
!!{proof}s ना लागू होते:
\[
\AxiomC{}
\RightLabel{$\pi_1'$}
\Deduce$\Gamma \fCenter \Delta', !B \land !C, !B, !A$
\AxiomC{}
\RightLabel{$\pi_3'$}
\Deduce$!A, \Gamma \fCenter \Delta', !B \land !C, !B$
\RightLabel{\CutCS}
\BinaryInf$\Gamma \fCenter \Delta', !B \land !C, !B$
\DisplayProof
\]
आणि
\[
\AxiomC{}
\RightLabel{$\pi_1''$}
\Deduce$\Gamma \fCenter \Delta', !B \land !C, !C, !A$
\AxiomC{}
\RightLabel{$\pi_4'$}
\Deduce$!A, \Gamma \fCenter \Delta', !B \land !C, !C$
\RightLabel{\CutCS}
\BinaryInf$\Gamma \fCenter \Delta', !B \land !C, !C$
\DisplayProof
\]
येथे \Log{G3c} मधील
!!{proof}s $\pi_1'$ आणि
$\pi_1''$ या $\pi_1$
पासून !!{height} न वाढवणाऱ्या
दुर्बलीकरणाने
(\olref[seq][adm]{prop:weak-G3c-adm})
मिळतात: एकदा उजवीकडे $!B$
आणि एकदा $!C$ घालून.
!!{proof}s $\pi_3'$ आणि
$\pi_4'$ या अनुक्रमे $\pi_3$
आणि $\pi_4$ पासून
!!{height} न वाढवणाऱ्या
दुर्बलीकरणाने मिळतात:
यात $\pi_3$ व $\pi_4$ यांची
उंची टिकवून उजवीकडे $!B \land !C$
घालून.
%READERNOTE{OLINC-342: तिसऱ्या वृक्षात स्रोताचे Delta हे Delta' हवे; अन्यथा संदर्भातील B-and-C ची अतिरिक्त प्रत येते.}

विगमन गृहीतकानुसार
\Log{G3c} मधील
$\Gamma \fCenter \Delta', !B \land !C, !B$
आणि
$\Gamma \fCenter \Delta', !B \land !C, !C$
यांच्या !!{proof}s~$\pi_5$ आणि~$\pi_6$
मिळतात. $\RightR{\land}$
नियमाने त्यांना जोडून
$\Gamma \Sequent \Delta',
!B \land !C, !B \land !C$
ची !!a{proof} मिळते:
\[
\AxiomC{}
\RightLabel{$\pi_5$}
\Deduce$\Gamma \fCenter \Delta', !B \land !C, !B$
\AxiomC{}
\RightLabel{$\pi_6$}
\Deduce$\Gamma \fCenter \Delta', !B \land !C, !C$
\RightLabel{\RightR{\land}}
\BinaryInf$\Gamma \fCenter \Delta', !B \land !C, !B \land !C$
\DisplayProof
\]
$\Gamma \Sequent \Delta$,
म्हणजे $\Gamma \fCenter
\Delta', !B \land !C$,
याची !!a{proof} मिळवण्यासाठी
\olref[seq][inv]{prop:G3c-cont-adm},
म्हणजे आकुंचनाची ग्राह्यता,
लावा.
%READERNOTE{OLINC-343: अंतिम वृक्षात Delta' आणि B-and-C च्या दोन प्रती हव्यात; स्रोताच्या Delta मुळे तीन प्रती येतात.}

\paragraph{प्रसंग E: कटची कोटी कमी करणे.}
शेवटचा प्रसंग असा की $\pi_1$
आणि $\pi_2$ दोन्हींचा शेवट
अशा अनुमानांनी होतो ज्यांत
$!A$ मुख्य असते. $!A$ च्या
संभाव्य !!{main operator} नुसार
सहा प्रसंग होतात.

प्रत्येक प्रसंगात कटचे
!!{formula}~$!A$ असलेले
\CutCS{} अनुमान आपण $!A$
च्या उप-!!{formula}s वरील
एक किंवा अधिक कटांत बदलतो;
म्हणजे ज्या अनुमानांत $!A$
हे मुख्य !!{formula} असते,
त्यांच्या सक्रिय !!{formula}s
वर कट करतो. म्हणून या
प्रसंगांना कट-अनुमानाचे
“न्यूनीकरण” किंवा “परिवर्तन”
म्हणतात.

\begin{enumerate}
\item $!A \ident \lnot !B$
असेल, तर $\pi$ चा शेवट असा:
\[
\AxiomC{}
\RightLabel{$\pi_1'$}
\Deduce$!B, \Gamma \fCenter \Delta$
\RightLabel{\RightR{\lnot}}
\UnaryInf$\Gamma \fCenter \Delta, \lnot !B$
\LeftSubproofLabel{$\pi_1$}
\AxiomC{}
\RightLabel{$\pi_2'$}
\Deduce$\Gamma \fCenter \Delta, !B$
\RightLabel{\LeftR{\lnot}}
\UnaryInf$\lnot !B, \Gamma \fCenter \Delta$
\RightSubproofLabel{$\pi_2$}
\RightLabel{\CutCS}
\BinaryInf$\Gamma \fCenter \Delta$
\DisplayProof
\]
विगमन गृहीतकानुसार पुढील
सिद्धतेतून \CutCS{} काढता येतो:
\[
\AxiomC{}
\RightLabel{$\pi_2'$}
\Deduce$\Gamma \fCenter \Delta, !B$
\AxiomC{}
\RightLabel{$\pi_1'$}
\Deduce$!B, \Gamma \fCenter \Delta$
\RightLabel{\CutCS}
\BinaryInf$\Gamma \fCenter \Delta$
\DisplayProof
\]

\item
$!A \ident (!B \lor !C)$ असेल,
तर $\pi$ चा शेवट असा आहे:
\[
\AxiomC{}
\RightLabel{$\pi_1'$}
\Deduce$\Gamma \fCenter \Delta, !B, !C$
\RightLabel{\RightR{\lor}}
\UnaryInf$\Gamma \fCenter \Delta, !B \lor !C$
\LeftSubproofLabel{$\pi_1$}
\AxiomC{}
\RightLabel{$\pi_2'$}
\Deduce$!B, \Gamma \fCenter \Delta$
\AxiomC{}
\RightLabel{$\pi_2''$}
\Deduce$!C, \Gamma \fCenter \Delta$
\RightLabel{\LeftR{\lor}}
\BinaryInf$!B \lor !C, \Gamma \fCenter \Delta$
\RightSubproofLabel{$\pi_2$}
\RightLabel{\CutCS}
\BinaryInf$\Gamma \fCenter \Delta$
\DisplayProof
\]
\CutCS{} ने संपणारी पुढील
!!{proof} विचारात घ्या:
\[
\AxiomC{}
\RightLabel{$\pi_1'$}
\Deduce$\Gamma \fCenter \Delta, !B, !C$
\AxiomC{}
\RightLabel{$\pi_2'''$}
\Deduce$!C, \Gamma \fCenter \Delta, !B$
\RightLabel{\CutCS}
\BinaryInf$\Gamma \fCenter \Delta, !B$
\DisplayProof
\]
येथे $\pi_2'''$ ही $\pi_2''$
पासून !!{height} न वाढवणारे
दुर्बलीकरण लावून मिळते.
$\depth{!C} < \depth{!B \lor !C}$
असल्यामुळे तिची कट-कोटी
$< \cutr{\pi}$ आहे.
विगमन गृहीतकानुसार
$\Gamma \fCenter \Delta, !B$
ची \Log{G3c} मधील
!!a{proof}~$\pi_1''$ मिळते.
आता \CutCS{} ने संपणारी
पुढील !!{proof} पाहा:
\[
\AxiomC{}
\RightLabel{$\pi_1''$}
\Deduce$\Gamma \fCenter \Delta, !B$
\AxiomC{}
\RightLabel{$\pi_2'$}
\Deduce$!B, \Gamma \fCenter \Delta$
\RightLabel{\CutCS}
\BinaryInf$\Gamma \fCenter \Delta$
\DisplayProof
\]
हिची कट-कोटी
$= \depth{!B} < \depth{!B \lor !C}$
असल्यामुळे विगमन गृहीतक
पुन्हा लागू होते. त्यातून
$\Gamma \fCenter \Delta$
ची !!a{proof} $\pi'$ मिळते.

\item $!A \ident (!B \land !C)$.
हा प्रसंग सरावासाठी आहे.

\item $!A \ident (!B \lif !C)$
असेल, तर $\pi$ चा शेवट असा:
\[
\AxiomC{}
\RightLabel{$\pi_1'$}
\Deduce$!B, \Gamma \fCenter \Delta, !C$
\RightLabel{\RightR{\lif}}
\UnaryInf$\Gamma \fCenter \Delta, !B \lif !C$
\LeftSubproofLabel{$\pi_1$}
\AxiomC{}
\RightLabel{$\pi_2'$}
\Deduce$\Gamma \fCenter \Delta, !B$
\AxiomC{}
\RightLabel{$\pi_2''$}
\Deduce$!C, \Gamma \fCenter \Delta$
\RightLabel{\LeftR{\lif}}
\BinaryInf$!B \lif !C, \Gamma \fCenter \Delta$
\RightSubproofLabel{$\pi_2$}
\RightLabel{\CutCS}
\BinaryInf$\Gamma \fCenter \Delta$
\DisplayProof
\]
\CutCS{} ने संपणाऱ्या
पुढील !!{proof} ची कट-कोटी
$\pi$ पेक्षा कमी आहे:
\[
\AxiomC{}
\RightLabel{$\pi_1'$}
\Deduce$!B, \Gamma \fCenter \Delta, !C$
\AxiomC{}
\RightLabel{$\pi_2''$}
\Deduce$!C, \Gamma \fCenter \Delta$
\RightLabel{\CutCS}
\BinaryInf$!B, \Gamma \fCenter \Delta$
\DisplayProof
\]
विगमन गृहीतकानुसार
$!B, \Gamma \fCenter \Delta$
ची !!a{proof}~$\pi_1''$ मिळते.
आता \CutCS{} ने संपणारी
पुढील !!{proof} पाहा:
\[
\AxiomC{}
\RightLabel{$\pi_2'$}
\Deduce$\Gamma \fCenter \Delta, !B$
\AxiomC{}
\RightLabel{$\pi_1''$}
\Deduce$!B, \Gamma \fCenter \Delta$
\RightLabel{\CutCS}
\BinaryInf$\Gamma \fCenter \Delta$
\DisplayProof
\]
तिची कट-कोटीही $\pi$
च्या कट-कोटीपेक्षा कमी आहे.
म्हणून विगमन गृहीतकानुसार
$\Gamma \fCenter \Delta$
ची \Log{G3c}-!!{proof}~$\pi'$
मिळते.
\item $!A \ident \lforall[x][!B(x)]$
असेल, तर $\pi$ चा शेवट असा:
\[
\AxiomC{}
\RightLabel{$\pi_1'$}
\Deduce$\Gamma \fCenter \Delta, !B(c)$
\RightLabel{\RightR{\lforall}}
\UnaryInf$\Gamma \fCenter \Delta, \lforall[x][!B(x)]$
\LeftSubproofLabel{$\pi_1$}
\AxiomC{}
\RightLabel{$\pi_2'$}
\Deduce$!B(t), \lforall[x][!B(x)], \Gamma \fCenter \Delta$
\RightLabel{\LeftR{\lforall}}
\UnaryInf$\lforall[x][!B(x)], \Gamma \fCenter \Delta$
\RightSubproofLabel{$\pi_2$}
\RightLabel{\CutCS}
\BinaryInf$\Gamma \fCenter \Delta$
\DisplayProof
\]
विगमन गृहीतकानुसार
पुढील सिद्धतेतून \CutCS{}
काढता येतो:
\[
\AxiomC{}
\RightLabel{$\pi_1''$}
\Deduce$!B(t), \Gamma \fCenter \Delta, \lforall[x][!B(x)]$
\AxiomC{}
\RightLabel{$\pi_2'$}
\Deduce$!B(t), \lforall[x][!B(x)], \Gamma \fCenter \Delta$
\RightLabel{\CutCS}
\BinaryInf$!B(t), \Gamma \fCenter \Delta$
\DisplayProof
\]
येथे $\pi_1''$ ही $\pi_1$
पासून (फक्त~$\pi_1'$ पासून
नव्हे!) !!{height} न वाढवणारे
दुर्बलीकरण लावून मिळते.
(या !!{proof} ची कट-कोटी
तेवढीच, पण कटची !!{height}
कमी आहे.) त्यामुळे
$!B(t), \Gamma \fCenter \Delta$
ची !!a{proof}~$\pi_2''$ मिळते.

!!{proof}~$\pi_1'$ ची
अंतिम क्रमवर्ती
$\Gamma \Sequent \Delta, !B(c)$
आहे, जेथे $c$ हा
$\RightR{\lforall}$ अनुमानाचा
आयगेन चर आहे. म्हणून $c$ हे
$!B(x)$, $\Gamma$ किंवा~$\Delta$
यांत येत नाही. !!{proof}~$\pi$
नियमित आहे असे मानू; त्यामुळे
$c$ हा फक्त पुढील
$\RightR{\lforall}$ अनुमानाचा
आयगेन चर आहे आणि फक्त त्याच्या
वर येतो. म्हणजे तो फक्त~$\pi_1'$
मध्ये येतो आणि $\pi_1'$ मधील
कोणत्याही अनुमानाचा आयगेन चर
नसतो. $c$ हे $\pi_2$ मध्ये
येत नसल्यामुळे $t$ मध्येही
येत नाही.
\olref[seq][qua]{lem:subst-regular}
नुसार !!{proof}~$\Subst{\pi_1'}{t}{c}$
ही $\Gamma \Sequent \Delta, !B(t)$
ची !!a{proof} आहे. आता पुढील
!!{proof} ला विगमन गृहीतक लावू:
\[
\AxiomC{}
\RightLabel{$\Subst{\pi_1'}{t}{c}$}
\Deduce$\Gamma \fCenter \Delta, !B(t)$
\AxiomC{}
\RightLabel{$\pi_2''$}
\Deduce$!B(t), \Gamma \fCenter \Delta$
\RightLabel{\CutCS}
\BinaryInf$\Gamma \fCenter \Delta$
\DisplayProof
\]
(कटचे !!{formula}~$!B(t)$
असल्यामुळे विगमन गृहीतक
लागू होते; या !!{proof} ची
कट-कोटी $\pi$ पेक्षा कमी आहे.)
\item $!A \ident \lexists[x][!B(x)]$
हा प्रसंग सरावासाठी ठेवतो.
\end{enumerate}
\end{proof}

\begin{prob}
\olref{lem:cut-adm-G3c} च्या
सिद्धतेतील प्रसंग~D चा
पुढील उपप्रसंग तपासा:
$\pi_1$ चा शेवट
$\LeftR{\lif}$ ने होतो,
पण त्याच्या अंतिम अनुमानात
$!A$ मुख्य नसते.
(टीप: येथे नेमके काय सिद्ध
करायचे, म्हणजे या प्रसंगात
$\pi$ चे रूप काय, हे स्पष्ट
मांडणेही सरावाचा भाग आहे.)
\end{prob}

\begin{prob}
\olref{lem:cut-adm-G3c} च्या
सिद्धतेतील प्रसंग~D चा
पुढील उपप्रसंग तपासा:
$\pi_1$ चा शेवट
$\LeftR{\lforall}$ ने होतो,
पण त्याच्या अंतिम अनुमानात
$!A$ मुख्य नसते.
\end{prob}

\begin{prob}
\olref{lem:cut-adm-G3c} च्या
सिद्धतेतील प्रसंग~E चा
पुढील उपप्रसंग तपासा:
$\pi_1$ आणि $\pi_2$
या दोन्हींच्या अंतिम अनुमानात
$!A$ मुख्य आहे आणि
$!A \ident (!B \land !C)$.
\end{prob}

\begin{prob}
\olref{lem:cut-adm-G3c} च्या
सिद्धतेतील प्रसंग~E चा
पुढील उपप्रसंग तपासा:
$\pi_1$ आणि $\pi_2$
या दोन्हींच्या अंतिम अनुमानात
$!A$ मुख्य आहे आणि
$!A \ident \lexists[x][!B(x)]$.
\end{prob}

प्रसंग~E मध्ये विगमन
गृहीतक ज्या \CutCS{} ने
संपणाऱ्या !!{proof} वर
लावतो, तिची कटची !!{height}
मूळ सिद्धतेच्या कटच्या
!!{height} पेक्षा मोठी
असू शकते. उदाहरणार्थ,
$!A \ident (!B \lif !C)$
असताना $\pi$ चे दोन
\CutCS{} अनुमाने असणाऱ्या
!!a{proof} मध्ये रूपांतर
केले:
\[
\AxiomC{}
\RightLabel{$\pi_2'$}
\Deduce$\Gamma \fCenter \Delta, !B$
\AxiomC{}
\RightLabel{$\pi_1'$}
\Deduce$!B, \Gamma \fCenter \Delta, !C$
\AxiomC{}
\RightLabel{$\pi_2''$}
\Deduce$!C, \Gamma \fCenter \Delta$
\RightLabel{\CutCS}
\BinaryInf$!B, \Gamma \fCenter \Delta$
\RightSubproofLabel{$\pi_1''$}
\RightLabel{\CutCS}
\BinaryInf$!B, \Gamma \fCenter \Delta$
\DisplayProof
\]
वरचा \CutCS{}, जो
$\pi_1'$ आणि $\pi_2''$
यांच्यामध्ये आहे, त्याची
कट-कोटी आणि कटची !!{height}
दोन्ही $\pi$ पेक्षा कमी आहेत.
पण विगमन गृहीतक लावल्याने
मिळालेली !!{proof}~$\pi_1''$
हिची कटची !!{height} आता
$\pi$ पेक्षा कमी असण्याची
गरज नाही. तरी खालच्या
\CutCS{} चे कटचे !!{formula}
$!B$ असल्यामुळे त्याची
\emph{कोटी} $\pi$ च्या
कट-कोटीपेक्षा कमी आहे.

\olref{lem:cut-adm-G3c} नुसार
!!a{proof} मधील एक सर्वोच्च
\CutCS{} अनुमान काढता येते.
हे उपप्रमेय सर्वोच्च~\CutCS{}
ने संपणाऱ्या उप-!!{proof}s
वर एकामागोमाग लावून
!!a{proof} मधील कितीही
\CutCS{} अनुमाने काढता
येतात.

\begin{thm}\ollabel{thm:cut-elim-G3c-topmost}
$\Log{G3c} + \CutCS$ मधील
कोणतीही !!{proof} $\pi$
ही~\Log{G3c} मधील
!!a{proof} मध्ये रूपांतरित
करता येते.
\end{thm}

\begin{proof}
  $\pi$ मधील \CutCS{}
  अनुमानांची संख्या~$n$
  हिच्यावर विगमन करू.
  $n=0$ असेल, तर काही
  करायचे नाही: $\pi$
  ही आधीच~\Log{G3c}
  मधील !!a{proof} आहे.
  अन्यथा, सर्वोच्च
  \CutCS{} अनुमानावर
  संपणारी उप-!!{proof}~$\pi'$
  विचारात घ्या. तिचे
  एकमेव \CutCS{} अंतिम
  अनुमान असते.
  \olref{lem:cut-adm-G3c}
  लावून तिचे \CutCS{}
  विरहित सिद्धता~$\pi''$
  मध्ये रूपांतर करा आणि
  मूळ सिद्धतेत
  उप-!!{proof}~$\pi'$
  च्या जागी~$\pi''$ ठेवा.
  मिळालेल्या !!a{proof}
  मध्ये $n-1$ \CutCS{}
  अनुमाने उरतात; म्हणून
  विगमन गृहीतक लागू होते.
\end{proof}

\end{document}
